両側の評価は一つの関数
\begin{equation}
\varphi(\rho,a,b)=\frac{\rho\,a}{1-\rho\,(1-a-b)}\qquad(a=0\text{では}\varphi=0)
\label{eq:phi}
\end{equation}
で書ける．$\varphi$は$\rho$と$a$について増加，$b$について減少し，漏れがなければ
\begin{equation*}
\varphi(1,a,b)=\frac{a}{a+b}
\end{equation*}
である．出口$a$と出口$b$の競争に，訪問ごとの生存率$\rho$が掛かる量と読める．

\appblock
\textbf{\TEscape（逃げる確率：未知の側）．}支持画像も種も含まない節点の集合$Q$について，$Q$を初めて出る先が$z$である割引確率を$H_Q(y,z)$とすると，
\begin{equation*}
h(y)=\sum_{z\in\partial Q}H_Q(y,z)\,h(z)\qquad\text{（等式）}
\end{equation*}
であり，種でない外へ逃げる確率
\begin{equation*}
\pi_Q(y)=\sum_{z\in\partial Q\setminus\hat{Q}}H_Q(y,z)
\end{equation*}
について$h(y)\le\pi_Q(y)$である．$y\in Q$の重みのうち種でない外へ向かう割合を$e_y$，種へ向かう割合を$\zeta_y$とすると，
\begin{equation*}
\pi_Q\le\varphi\big(\max_Q\rho,\ \max_Qe,\ \min_Q\zeta\big).
\end{equation*}

\appblock
\textbf{\TCapture（捕まる確率：IDの側）．}支持画像も種も含まない節点の集合$B$について，$y\in B$の重みのうち支持画像へ向かう割合を$\omega_y$，$B$にも支持画像にも属さない節点へ向かう割合を$\eta_y$とすると，$B$を出ずに支持画像に着く確率$\kappa_B$について
\begin{equation*}
h\ge\kappa_B\ge\varphi\big(\min_B\rho,\ \min_B\omega,\ \max_B\eta\big)
\end{equation*}
である．また，すべてのラベルなし画像で
\begin{equation*}
s(x)\ge\rho_x\,\omega_x
\end{equation*}
である．

\appblock
\textbf{\CSep（分離）．}種でないテスト画像$y$が\TEscape{}の$Q$に属し，各較正画像$x'$について$\ell(x')\le s(x')$（例えば$\ell(x')=\rho_{x'}\omega_{x'}$）なら，
\begin{equation}
p(y)\le\frac{1+|\{x'\in\mathcal{C}:\ell(x')\le\pi_Q(y)\}|}{|\mathcal{C}|+1}.
\label{eq:sep}
\end{equation}
特に$\pi_Q(y)<\min_{x'\in\mathcal{C}}\ell(x')$なら$p(y)$は床$1/(|\mathcal{C}|+1)$に落ち，$y$は$\alpha\ge1/(|\mathcal{C}|+1)$のすべての水準でOODと判定される（証明は付録\ref{sec:app_h19}）．

\appblock
\textbf{読み方．}未知の塊では「種でない外への辺」と「種への辺」が競い，IDの近傍では「支持画像への辺」と「漏れる辺」が競う．同じ$\varphi$が両側を抑え，その大小が較正順位の分離\eqnref{sep}を決める．有効性（\TValid）と検知力（\TEscape{}--\CSep）は同じ量$s$についての主張である．種を増やすと$h$は単調に減り，塊のどの節点でも重みの割合$\zeta$以上が種へ向かえば，漏れがなくても逃げる確率は
\begin{equation*}
\pi_Q\le\frac{\max_Qe}{\max_Qe+\zeta}
\end{equation*}
を満たす．同じ未知クラスの画像が複数枚あるとき，種になった画像は残りの画像の重みの一部を占め，その逃げる確率を下げる．これが自己ラベル付けの効く理由である．
